\section{A universal polynomial quotient at 289}
\label{numerator:section}

The distinction between finite and infinite digit expansions has a second
algebraic aspect. For nonzero polynomials $P,Q\in\Z[t]$ with coefficients
in $\{0,1\}$, a representation $M=P(3)/Q(3)$ is called
\emph{universal} if $Q$ divides $P$, and \emph{local} otherwise.
This is the terminology of~\cite[Definition~1.1]{ADRS}, after removing
any common initial power of $t$. Indeed, for $M\equiv1\pmod3$ the
least nonzero digits of $P(3)$ and $Q(3)$ occur at the same position;
removing these equal powers of $t$ gives constant terms~1, without
changing polynomial divisibility. Since $Q$ is monic, divisibility
in $\Q[t]$ is equivalent to divisibility in $\Z[t]$.
The distinction concerns each representation, rather than just its
integer value.

The catalogue in~\cite[Section~4, pp.~24--25]{ADRS}, specifically
arXiv version~1, leaves open the existence of a local indecomposable
representation of~289. We answer this question by proving a stronger
statement about all finite representations, without an indecomposability
restriction.

\begin{theorem}[A common polynomial quotient for every representation]
\label{numerator:main}
Put
\[
 h(t)=1-t^2-t^3+t^4+t^5.
\]
For all nonzero polynomials $P,Q\in\Z[t]$ with coefficients in
$\{0,1\}$,
\begin{equation}\label{numerator:eq:universal}
 P(3)=289Q(3)\quad\Longrightarrow\quad P(t)=h(t)Q(t).
\end{equation}
In particular, every finite representation of~289 is universal, and
there is no local representation of~289.
\end{theorem}

\subsection{Returning carry paths and polynomial potentials}

We give the finite certificate and its all-length interpretation
explicitly. Let $M>1$ and $M\equiv1\pmod3$. An input digit $a$ and an
output digit $b$ in $\{0,1\}$ give the carry transition
\begin{equation}\label{numerator:eq:carry}
 3c'=c+Ma-b,\qquad c_0=0.
\end{equation}
Along a path, induction gives
\begin{equation}\label{numerator:eq:invariant}
 3^j c_j=M\sum_{i<j}a_i3^i-\sum_{i<j}b_i3^i.
\end{equation}
All reachable carries are nonnegative: the only negative numerator
possible from $c\ge0$ is $c=a=0,b=1$, and it is not divisible by~3.
The invariant also gives $c_j<M/2$, since
$\sum_{i<j}a_i3^i\le(3^j-1)/2$.
For each nonnegative carry all integral transitions are exactly
\begin{equation}\label{numerator:eq:all-edges}
\begin{array}{c|l}
 c\pmod3 & (a,b,c')\\\hline
 0 &(0,0,c/3),\quad(1,1,(c+M-1)/3),\\
 1 &(0,1,(c-1)/3),\\
 2 &(1,0,(c+M)/3).
\end{array}
\end{equation}
These are obtained by checking the four digit pairs modulo~3.

Every finite pair $P(3)=M Q(3)$, padded with high zero digits to a
common length $n$, determines a path from~0 to~0: its carries are
integers because each prefix difference is divisible by $3^j$.
Conversely, a path from~0 to~0 gives this integer equality by
\eqref{numerator:eq:invariant}. At least one input digit is~1 exactly
when the witness is positive.

\begin{lemma}[Polynomial potential certificate]
\label{numerator:potential}
Suppose $S$ contains every state on every path from~0 back to~0.
Let $h\in\Z[t]$ satisfy $h(3)=M$. Suppose polynomials
$H_c\in\Z[t]$, for $c\in S$, satisfy $H_0=0$ and
\begin{equation}\label{numerator:eq:potential}
 tH_{c'}=H_c+a h-b
\end{equation}
on every transition $c\xrightarrow{a,b}c'$ inside $S$.
Then every finite witness satisfies $P=hQ$.
\end{lemma}
\begin{proof}
Multiply \eqref{numerator:eq:potential} on the $j$th edge by $t^j$
and sum along a returning path of length~$n$. The potentials telescope:
\[
 t^nH_{c_n}-H_{c_0}
   =h\sum_{j<n}a_jt^j-\sum_{j<n}b_jt^j.
\]
The left side is zero. This is an identity in $\Z[t]$ for each finite
$n$, with no bound on the witness length.
\end{proof}

\subsection{The complete certificate for 289}

For $M=289$, define the disjoint sets
\begin{align*}
 S={}&\{0,1,4,11,12,32,33,37,44,48,96,100,107,111,112,\\
     &\hspace{18mm}132,133,140,143,144\},\\
 D={}&\{5,14,15,16,43,46,98,101,128,129,130,139\}.
\end{align*}
Table~\ref{numerator:tab:certificate} lists all transitions out of
$S\cup D$, using \eqref{numerator:eq:all-edges}, and gives the
polynomials $H_c$ on~$S$. A dash marks a state in~$D$.
Every listed destination is in $S\cup D$, so this set contains all
states reachable from~0. Every transition out of $D$ stays in $D$,
and $0\notin D$. Consequently a returning path cannot enter $D$;
it stays in $S$.

\begin{table}[tbp]
\centering
\small
\setlength{\tabcolsep}{7pt}
\begin{tabular}{rll}
\toprule
Carry $c$ & Potential $H_c(t)$ & All transitions $(a,b):c'$\\
\midrule
0 & $0$ & $(0,0):0;\quad (1,1):96$ \\
1 & $1$ & $(0,1):0$ \\
4 & $1+t$ & $(0,1):1$ \\
5 & $\text{--}$ & $(1,0):98$ \\
11 & $-1+t+t^{2}$ & $(1,0):100$ \\
12 & $t+t^{2}$ & $(0,0):4;\quad (1,1):100$ \\
14 & $\text{--}$ & $(1,0):101$ \\
15 & $\text{--}$ & $(0,0):5;\quad (1,1):101$ \\
16 & $\text{--}$ & $(0,1):5$ \\
32 & $-1-t+t^{2}+t^{3}$ & $(1,0):107$ \\
33 & $-t+t^{2}+t^{3}$ & $(0,0):11;\quad (1,1):107$ \\
37 & $1+t^{2}+t^{3}$ & $(0,1):12$ \\
43 & $\text{--}$ & $(0,1):14$ \\
44 & $-1+2t^{2}+t^{3}$ & $(1,0):111$ \\
46 & $\text{--}$ & $(0,1):15$ \\
48 & $t+2t^{2}+t^{3}$ & $(0,0):16;\quad (1,1):112$ \\
96 & $-t-t^{2}+t^{3}+t^{4}$ & $(0,0):32;\quad (1,1):128$ \\
98 & $\text{--}$ & $(1,0):129$ \\
100 & $1-t^{2}+t^{3}+t^{4}$ & $(0,1):33$ \\
101 & $\text{--}$ & $(1,0):130$ \\
107 & $-1+t^{3}+t^{4}$ & $(1,0):132$ \\
111 & $t+t^{3}+t^{4}$ & $(0,0):37;\quad (1,1):133$ \\
112 & $1+t+t^{3}+t^{4}$ & $(0,1):37$ \\
128 & $\text{--}$ & $(1,0):139$ \\
129 & $\text{--}$ & $(0,0):43;\quad (1,1):139$ \\
130 & $\text{--}$ & $(0,1):43$ \\
132 & $-t+2t^{3}+t^{4}$ & $(0,0):44;\quad (1,1):140$ \\
133 & $1-t+2t^{3}+t^{4}$ & $(0,1):44$ \\
139 & $\text{--}$ & $(0,1):46$ \\
140 & $-1-t+t^{2}+2t^{3}+t^{4}$ & $(1,0):143$ \\
143 & $-1+t^{2}+2t^{3}+t^{4}$ & $(1,0):144$ \\
144 & $t^{2}+2t^{3}+t^{4}$ & $(0,0):48;\quad (1,1):144$ \\
\bottomrule
\end{tabular}
\caption{A closed carry certificate for $M=289$. The 32 rows list
all 42 edges out of the displayed states. Potentials are assigned
exactly on $S$; the other 12 states form the closed set $D$.}
\label{numerator:tab:certificate}
\end{table}

For each of the 26 edges with both endpoints in $S$, the displayed
polynomials satisfy \eqref{numerator:eq:potential}. This is checked
by subtracting the two sides coefficient by coefficient, with
$h=1-t^2-t^3+t^4+t^5$. For example, the edges from~0, 4 and~144 give
\[
 tH_{96}=h-1,\qquad tH_1=H_4-1,\qquad
 tH_{144}=H_{144}+h-1.
\]
The table is therefore a finite, explicit certificate for every
hypothesis of Lemma~\ref{numerator:potential}. Merely checking
$H_c(3)=c$ would not suffice; the full polynomial identities are the
input to the lemma.

\begin{proof}[Proof of Theorem~\ref{numerator:main}]
We have $h(3)=289$ and $H_0=0$. The preceding closure checks show that
all returning paths stay in~$S$, and the table verifies the potential
identity on every edge they can use. Apply
Lemma~\ref{numerator:potential}.
\end{proof}

The theorem is not vacuous. Direct multiplication gives
\begin{equation}\label{numerator:eq:witness}
 (1-t^2-t^3+t^4+t^5)(1+t^2+t^3+t^5)
       =1+t^7+t^9+t^{10},
\end{equation}
which at $t=3$ is $289\cdot280=80920$.
The accompanying verifier independently constructs all allowable
transitions in the a priori range $0\le c\le144$, checks closure and
all polynomial coefficients, and checks \eqref{numerator:eq:witness}.
Its tests of additional long returning paths are regression tests;
the all-length conclusion follows from the displayed certificate and
telescoping argument.

\subsection{An explicit family of local finite representations}

The same polynomial viewpoint produces a family for which a specified
representation is local. To distinguish it from the sporadic family
$M_k$ in Theorem~\ref{thm:family}, write
\[
 N_k=4\cdot3^k-2\qquad(k\ge1).
\]

\begin{theorem}[Explicit finite witnesses]
\label{numerator:local-family}
For every $k\ge1$, $N_k\in\mathcal V$. There are $0,1$ polynomials
$P_k,Q_k$ with
\[
 P_k(3)=N_k Q_k(3),\qquad
 \deg Q_k=k^2-1,\qquad\deg P_k=k(k+1).
\]
For $k\ge2$, the representation specified below is local.
\end{theorem}
\begin{proof}
Set
\begin{equation}\label{numerator:eq:q-family}
 Q_k(t)=\sum_{j=0}^{k-1}t^{(k+1)j}
                         \sum_{i=0}^{k-j-1}t^i.
\end{equation}
The exponent intervals in the outer sum are disjoint, and the final
one ends at $(k+1)(k-1)=k^2-1$. Thus $Q_k$ is a $0,1$ polynomial
of the asserted degree, with constant term~1.
For $k=1$, take $P_1=1+t^2$ and $Q_1=1$.
For $k\ge2$, put $n=k(k+1)$ and set
\begin{equation}\label{numerator:eq:p-family}
 P_k(t)=1+\sum_{i=2}^{k}t^{ki}
       +\sum_{j=1}^{k-2}t^{k+(k+1)j}+t^{n-1}+t^n.
\end{equation}
The exponents in the second sum have distinct nonzero residues
$1,\ldots,k-2$ modulo~$k$; the exponents in the first sum are
multiples of~$k$. The final two exponents exceed all preceding ones.
Hence all displayed exponents are distinct, so $P_k$ is also a
$0,1$ polynomial, of degree~$n$.

Write $S_k(t)=\sum_{i=0}^{k-1}t^i$. Shifting the blocks in
\eqref{numerator:eq:q-family} and subtracting successive blocks gives
\begin{equation}\label{numerator:eq:telescoping-one}
 (t^{k+1}-1)Q_k=\sum_{i=2}^{k+1}t^{ki}-S_k.
\end{equation}
A second identity is
\begin{equation}\label{numerator:eq:telescoping-two}
 (t^k-1)Q_k=\sum_{j=0}^{k-1}t^{k+(k+1)j}-S_k.
\end{equation}
For completeness, multiplying \eqref{numerator:eq:q-family} by $t-1$
gives
\[
 (t-1)Q_k=\sum_{i=1}^{k}t^{ki}
                     -\sum_{j=0}^{k-1}t^{(k+1)j}.
\]
Multiplication by $t^k-1$, followed by
\[
 (t^k-1)\sum_{i=1}^{k}t^{ki}=t^{k(k+1)}-t^k,
 \qquad
 (t^{k+1}-1)\sum_{j=0}^{k-1}t^{(k+1)j}=t^{k(k+1)}-1,
\]
proves \eqref{numerator:eq:telescoping-two} after cancellation of the
nonzero polynomial $t-1$.
Adding \eqref{numerator:eq:telescoping-one} and
\eqref{numerator:eq:telescoping-two}, and separating their extreme
terms, now gives
\begin{equation}\label{numerator:eq:division}
 \begin{split}
 P_k&=(t^{k+1}+t^k-2)Q_k+R_k,\\
 R_k&=3+2\sum_{i=1}^{k-1}t^i-t^k.
 \end{split}
\end{equation}
At $t=3$, the remainder is zero, so
$P_k(3)=(4\cdot3^k-2)Q_k(3)$.
But $R_k$ is a nonzero polynomial with
$\deg R_k=k<k^2-1=\deg Q_k$ for every $k\ge2$.
Thus \eqref{numerator:eq:division} is polynomial division with a
nonzero remainder. The displayed representation is local.
\end{proof}

\begin{remark}[Scope of the two conclusions]
Theorem~\ref{numerator:main} determines the type of \emph{every}
representation of one integer. Theorem~\ref{numerator:local-family}
constructs \emph{one} representation for each member of an infinite
family; it does not exclude other universal representations or claim
minimal witness length. The output length is $k(k+1)+1$, hence
$O((\log N_k)^2)$ on this family. Since the family has only
$O(\log X)$ members up to~$X$, this construction does not improve
the counting exponent of Theorem~\ref{ep:main}. Neither conclusion
proves the finiteness of good powers of two.
\end{remark}
